\section{Sign changes of the discrete law}\label{sec:signs}

The density $\La$ has $1$ or $3$ critical points when $b=1$. The discrete law can have
more local extrema, and the extra ones sit at the ends. In this section we count the
sign changes of $k\mapsto P_N(k+1)-P_N(k)$, which are the strict local extrema of $P_N$
when no two neighbouring values are equal. For a finite or infinite real sequence
$(s_i)$ we write $V(s)$ for the number of sign changes of the sequence of its nonzero
terms.

\begin{introtheorem}[a lower bound]\label{thm:signs}
Assume (RV$_\alpha$), (SPAN) and (SRT$_\pm$), and let the initial run be such that
$NP_N(k)=\La(k/N)+o(1)$ uniformly on compact subsets of $(0,1)$ (for instance a fresh
start, or under (LB$_\pm$) any initial run with $n\nu_\pm(n)\to0$) and such that
$W^\pm_N=o(T_\pm(N))$. Assume that the limits
\[
 \kappa_\plus(0)=\lim_N\frac{Q_\plus(N)}{T_\plus(N)},\qquad \kappa_\minus(0)=\lim_N\frac{Q_\minus(N)}{T_\minus(N)}
\]
exist and are finite, and let $\kappa_\pm(j)$ for $j\ge1$ be the end profiles
of~\eqref{eq:kappa-prep} and its mirror. Assume that $\kappa_\pm(j+1)\ne\kappa_\pm(j)$ for
all $j\ge0$ and that $\kappa_\pm$ is strictly decreasing from some point on. Let
$\sigma_\pm=V\bigl((\kappa_\pm(j+1)-\kappa_\pm(j))_{j\ge0}\bigr)$, and assume that $\La$ has
finitely many critical points, $\sigma_L$ in number, each a strict local extremum. Then for
all large $N$,
\[
 V\bigl((P_N(k+1)-P_N(k))_{0\le k<N}\bigr)\ge\sigma_L+\sigma_\plus+\sigma_\minus .
\]
\end{introtheorem}

In the definition of $\kappa_\pm(0)$ we use the atom. By Proposition~\ref{prop:first-run},
$P_N(N)=Q_\plus(N)$ exactly, while the formula~\eqref{eq:kappa-prep} gives $0$ at $j=0$. If
$Q_\plus(N)/T_\plus(N)$ does not converge, the value at the end point can break the
pattern of the profile, and the theorem says nothing.

\begin{proof}
We use the following fact. If $(s_i)$ is a sequence and $(s_{i_l})$ a subsequence, then
the sequence of differences of $(s_i)$ has at least as many sign changes as the sequence
of differences of $(s_{i_l})$. Indeed, each nonzero difference $s_{i_{l+1}}-s_{i_l}$ is a
sum of consecutive differences of $(s_i)$, one of which has its sign, and choosing one
such difference for each $l$ gives a subsequence of the differences of $(s_i)$ with the
same sign pattern.

Let $c_1<\dots<c_s$ be the critical points of $\La$, $s=\sigma_L$. Since $\La\to\infty$ at
both ends, $\La$ decreases strictly on $(0,c_1]$ and increases strictly on $[c_s,1)$, and
the critical points alternate between strict minima and strict maxima, starting and ending
with a minimum. Choose $z_0\in(0,c_1)$ and $z_{s+1}\in(c_s,1)$. Then the differences of
$\La(z_0),\La(c_1),\dots,\La(c_s),\La(z_{s+1})$ have the signs $-,+,-,\dots,+$, with $s$
sign changes. Choose $J$ such that $\kappa_\plus$ and $\kappa_\minus$ are strictly
decreasing on $[J-1,\infty)$. Since a strictly decreasing nonnegative sequence is positive,
$\kappa_\pm(J)>0$.

Consider the subsequence of $(P_N(k))_{0\le k\le N}$ at the points
\[
 0,1,\dots,J,\ \lfloor z_0N\rfloor,\lfloor c_1N\rfloor,\dots,\lfloor c_sN\rfloor,\lfloor z_{s+1}N\rfloor,\ N-J,\dots,N .
\]
By Theorem~\ref{thm:initial}(f), its mirror, and the definition of $\kappa_\pm(0)$,
$P_N(N-j)/T_\plus(N)\to\kappa_\plus(j)$ and $P_N(j)/T_\minus(N)\to\kappa_\minus(j)$ for
$0\le j\le J$. Since consecutive values of $\kappa_\pm$ differ, for large $N$ the
differences at the two ends have the signs of the differences of $\kappa_\pm$. These have
$\sigma_\minus$ sign changes at the minus end and $\sigma_\plus$ at the plus end, all at
indices below $J$, and the differences next to $j=J$ have the sign of
$\kappa_\pm(J)-\kappa_\pm(J-1)<0$. In the bulk, $NP_N(\lfloor zN\rfloor)\to\La(z)$ at each
chosen point, so for large $N$ the bulk differences have the signs $-,+,\dots,+$ with $s$
sign changes. At the two junctions, $P_N(J)\sim T_\minus(N)\kappa_\minus(J)$ and
$P_N(N-J)\sim T_\plus(N)\kappa_\plus(J)$, while the bulk values are $O(1/N)$. Since
$NT_\pm(N)\to\infty$, the difference from $P_N(J)$ to $P_N(\lfloor z_0N\rfloor)$ is negative
and the difference from $P_N(\lfloor z_{s+1}N\rfloor)$ to $P_N(N-J)$ is positive. Reading
from left to right, the last difference at the minus end is negative, the junction is
negative and the first bulk difference is negative. The last bulk difference is positive,
the junction is positive, and the first difference at the plus end, which is
$P_N(N-J+1)-P_N(N-J)\sim T_\plus(N)(\kappa_\plus(J-1)-\kappa_\plus(J))$, is positive. So the
junctions add no sign change, and the subsequence has $\sigma_\minus+s+\sigma_\plus$ sign
changes. The fact above finishes the proof.
\end{proof}

The hypothesis that $\kappa_\pm$ is eventually strictly decreasing holds in the main
examples. If $f_\minus$ is strictly log-convex, then by Theorem~1 of de Bruijn and
Erd\H{o}s~\cite{debruijnerdos1953} the renewal sequence $u_\minus$ is strictly log-convex.
Since $u_\minus(j)\to0$ under (SRT$_\minus$), the ratios $u_\minus(j+1)/u_\minus(j)$, which
increase, stay below $1$, so $u_\minus$ is strictly decreasing. The zeta law
$f(n)=n^{-\alpha-1}/\zeta(\alpha+1)$ is strictly log-convex.

\begin{example}[the fair start and the fresh start]\label{ex:counts}
Let both run laws be the zeta law, so $b=1$, and let the start be fair. Then
$Q_\pm(N)=\frac12T(N)$, so $\kappa_\pm(0)=\frac12$, and~\eqref{eq:kappa-prep} gives
$\kappa_\pm(j)=\frac12u(j)+\frac12\sum_{d\le j}f(d)u(j-d)=u(j)$ for $j\ge1$. Since $u$ is
strictly decreasing, $\sigma_\pm=1$ if $u(1)=f(1)>\frac12$ and $\sigma_\pm=0$ if $f(1)<\frac12$.
Now $f(1)=1/\zeta(1+\alpha)$, and $\zeta(1.7)=2.054\ldots$ and $\zeta(1.9)=1.750\ldots$. So
Theorem~\ref{thm:signs} and Theorem~\ref{thm:shape} give at least $1$ sign change for
$\alpha=0.3$ and $\alpha=0.5$, at least $3$ for $\alpha=0.7$, and at least $3+1+1=5$ for
$\alpha=0.9$. For a fresh plus start with different laws, $\kappa_\plus=(1,u_\minus(1),\dots)$
is strictly decreasing and $\kappa_\minus=(0,u_\plus(1),u_\plus(2),\dots)$ has one sign
change, so the bound is $\sigma_L+0+1$.
\end{example}

In every case of Section~\ref{sec:verify} the computed number of sign changes equals the
lower bound. We expect this in general.

\begin{conjecture}[the exact count]\label{conj:count}
Under the hypotheses of Theorem~\ref{thm:signs}, (LB$_\pm$), the condition
$|f_\pm(n)-f_\pm(n+1)|=O(f_\pm(n)/n)$ and strictly log-convex run laws, the number of sign
changes of $k\mapsto P_N(k+1)-P_N(k)$ equals $\sigma_L+\sigma_\plus+\sigma_\minus$ for all
large $N$.
\end{conjecture}

Two lemmas would prove it. The first is a difference form of Theorem~\ref{thm:local},
$N^2[P_N(k+1)-P_N(k)]\to\La'(x)$ uniformly for $k/N\to x$ in compact subsets of $(0,1)$.
The second is a difference form of Theorem~\ref{thm:ends}(b): for $J\le j\le\varepsilon N$
the sign of $P_N(N-j-1)-P_N(N-j)$ is that of $u_\minus(j+1)-u_\minus(j)<0$. The second
needs a difference form of the strong renewal theorem,
$u(j)-u(j+1)\sim(1-\alpha)u(j)/j$, and we do not know whether it holds under these
hypotheses when $\alpha\le1/2$.
